\section{The sharp cubic list bound on fourteen coordinates}
\label{app:global-seven-cubics}

\begin{theorem}\label{thm:global-seven-cubics}
There is an absolute constant $p_0$ with the following property.  If $k$
has characteristic zero or characteristic greater than $p_0$, then for
every fourteen-element set $\Omega\subset k$ and every word on $\Omega$,
at most seven distinct polynomials of degree at most three have at least
seven agreements with the word.  The bound is attained in characteristic
zero and over arbitrarily large prime fields.
\end{theorem}

The lower bound is Theorem~\ref{thm:seven-cubic-bank} with $m=1$.
The upper bound proved here is specific to cubics on fourteen coordinates;
it does not assert a universal list bound of seven for the quarter-rate
length-twenty-eight pullback.  The proof below includes two finite exact
certificates: an exhaustive incidence enumeration and a rational
polynomial-identity certificate.  Their algorithms and the mathematical
conditions they certify are specified explicitly.

\begin{proof}
\smallskip\noindent\textbf{Saturation and incidence reduction.}
Suppose eight distinct cubics have at least seven agreements.  For each
coordinate let $s_x$ be their matching multiplicity.  Pairwise root
counting and convexity give
\[
 \sum_xs_x\geq56,\qquad
 84\leq\sum_x\binom{s_x}{2}\leq3\binom82=84.
\]
Thus every coordinate has multiplicity four, every candidate has exactly
seven matches, and every pair has exactly three common matches.  The
fourteen matching sets form a $2$-$(8,4,3)$ design.
Its blocks are distinct: fixing a block $B$, the intersections $j$ with
the other thirteen blocks satisfy $\sum j=24$, $\sum\binom j2=12$, and
hence $\sum(j-1)(j-2)/2=1$.  All summands are nonnegative, whereas a
repeated block would contribute three.

Delete a candidate.  The seven coordinates at which it matched become
triple columns $T$; the other seven remain quadruple columns.  Write $C$
for the complementary triples of these quadruples in the remaining
seven candidates.  Both $T$ and $C$ are three-regular collections of
seven triples, and every pair has combined multiplicity two in $T,C$.
The extension back to eight candidates is forced: append index eight to
each $T$-triple and leave the complementary-$C$ quadruples unchanged.

\smallskip\noindent\textbf{The finite incidence certificate.}
Let
\[
 \mathcal F=(123,145,167,246,257,347,356),
\]
where strings denote subsets.  Every $2$-$(8,4,3)$ design arising above
has a deletion which, after relabeling, has $T=\mathcal F$ and one of
the following four ordered choices for $C$:
\[
\begin{array}{c|l}
4&123,145,167,246,257,347,356\\
5&123,145,167,247,256,346,357\\
6&123,146,157,247,256,345,367\\
7&124,136,157,237,256,345,467.
\end{array}                                                    \tag{G1}
\]
Here is the exhaustive algorithm verifying this finite assertion.
Enumerate nondecreasing lists of seven of the $35$ triples on seven
vertices, retaining exactly those with vertex degrees three and pair
degrees at most two.  Repetitions are permitted.  Encode the $21$ pair
degrees in base three.  Two lists are compatible exactly when their
digits sum to two, so a complementary-key lookup lists every ordered
pair $(T,C)$.  Quotient by simultaneous $S_7$ relabeling, without
interchanging $T,C$.  There are $12{,}780$ labelled halves, $6{,}150$
compatible ordered pairs, and eight ordered orbits.  For each
representative, form the forced eight-candidate extension and try all
eight deletions.  Each has a deletion of one of the four types (G1).
The other four types have deletion sets $\{0,3,6\}$, $\{1,5\}$,
$\{2,7\}$, and $\{0,3,6\}$; the four displayed types have deletion
sets $\{4\}$, $\{1,5\}$, $\{0,3,6\}$, and $\{2,7\}$, respectively.

The implementations \texttt{designs.cpp} and
\texttt{independent\_design\_set.py} independently verify the literal
set of all $6{,}150$ pairs, not only its cardinality.
\texttt{eight\_bank\_removal\_types.py} checks the deletion assertion.
All these files and the explicit representatives are in the accompanying
\texttt{fano\_seven} certificate directory.  The only enumeration
pruning discards partial lists whose vertex or pair degrees already
exceed the prescribed bounds, or whose vertices cannot reach degree
three.  Consequently no admissible incidence pattern is discarded.

We now exclude types 4--6 in characteristic different from two and
classify type 7 in characteristic zero.  We work over an algebraic
closure and with binary cubics, so projective coordinate changes are
permitted.  Every pair already agrees at three prescribed distinct
points.  Therefore an additional pair equality is impossible, including
an outsider sharing a quadruple's value.

\smallskip\noindent\textbf{Type 4.}
Normalize the seven quadruple coordinates, in the order of $C$, to
$\infty,0,1,u,v,w,z$.  All finite parameters are distinct and different
from zero and one.  On each $T$-triple, divide the three difference
identities by the common triple-coordinate factor.  The resulting three
binary quadratics, whose roots are the prescribed quadruple coordinates,
are linearly dependent.  For triples $167,246,347$ their determinants are
\[
 E_2=-uz+vw,\quad E_3=vw-vz-v+z,\quad E_5=uz-vz+v-z.
\]
All vanish, but $E_3-E_5-E_2=2(z-v)$ contradicts distinctness.

\smallskip\noindent\textbf{Type 5.}
Use the same normalization.  The corresponding quadratic determinants
for $123,145,167$ give
\[
 -uvw-uvz+uv+uwz+vwz-wz=0,
 \quad -uw+u+vz-v+w-z=0,\quad uz-vw=0.
\]
Substitution of $u=vw/z$ into the first two forces
$v(w+z-1)=z$.  With $s=w+z-1$ we obtain
$s\ne0$, $u=w/s$, $v=z/s$.
Subtract $P_1$ so that $P_1=0$.  All quadruple-compatible families are
\[
\begin{array}{ll}
 P_2=(X-w)(X-z)L_0,&P_3=s^2(X-u)(X-v)L_0,\\
 P_4=(X-z)(X-v)L_1,&P_5=(X-w)(X-u)L_1,\\
 P_6=(X-z)(X-u)L_2,&P_7=(X-w)(X-v)L_2,
\end{array}
\]
where $L_0=pX+q$, $L_1=\ell(X-1)-s(p+q)$, and
$L_2=\ell X+sq$.  Indeed the two forced quadruple roots determine each
quadratic factor, and the equalities at $0,1,\infty$ determine the
remaining linear factors as displayed.  In particular $\ell\ne0$ and
$p-\ell\ne0$, by the outsider guards at infinity.
Define
\begin{align*}
 A&=\ell z+pzs+qs(s+1),& B&=\ell w+pws+qs(s+1),\\
 C_z&=\ell(z-1)+s(w-1)p+s(s-1)q,
 &C_w&=\ell(w-1)+s(z-1)p+s(s-1)q,\\
 F_\pm&=\pm(z-w)\ell+s(s+1)p+2s^2q.
\end{align*}
The factors $A,B,C_z,C_w$ are nonzero: respectively they are, up to
nonzero scalar factors, $(P_2-P_4)(0)$, $(P_2-P_5)(0)$,
$(P_2-P_6)(1)$, and $(P_2-P_7)(1)$.
For triple $246$, divide the differences by their common factor $X-z$.
The resultant of
$s[(X-w)L_0-(X-v)L_1]$ and
$s[(X-w)L_0-(X-u)L_2]$ equals
$-(p-\ell)sAC_zF_+$.  The analogous resultant for triple $257$, after
dividing by $X-w$, is $-(p-\ell)sBC_wF_-$.
The required additional common roots force both resultants to vanish,
so $F_+=F_-=0$.  Their difference $2(z-w)\ell$ is nonzero, a contradiction.

\smallskip\noindent\textbf{Type 6.}
Relabel candidates by $(1,2,3,4,5,6,7)\mapsto(1,2,3,4,7,5,6)$.
Then $C=\mathcal F$ and
$T=(123,147,156,245,267,346,357)$.
Normalize the quadruple coordinates to $\infty,0,1,u,v,w,z$ and
subtract $P_1$.  Scale the nonzero quadruple value at zero to one, and
write the values at one and infinity as $K,L$, both nonzero.  Put
\[
\begin{array}{lll}
 Q_2=(X-w)(X-z),&Q_3=(X-u)(X-v),&Q_4=(X-v)(X-z),\\
 Q_5=(X-u)(X-w),&Q_6=(X-v)(X-w),&Q_7=(X-u)(X-z).
\end{array}
\]
The complete family is
\begin{align*}
 P_i&=Q_i[(K/Q_i(1)-1/Q_i(0))X+1/Q_i(0)] &&(i=2,3),\\
 P_i&=Q_i[L(X-1)+K/Q_i(1)] &&(i=4,5),\\
 P_i&=Q_i[LX+1/Q_i(0)] &&(i=6,7).
\end{align*}
The triple $123$ and the two other triples containing index one force,
with $s=w+z-1$,
\[
 v=\frac{wz(1-u)}{wz-us},\qquad
 L=\frac{K}{(1-v)(1-z)}-\frac1{uz}
  =\frac{K}{(1-u)(1-w)}-\frac1{vw}.             \tag{G2}
\]
The first denominator cannot vanish, since then its defining equation
would give $u=1$.  Substitution factors the difference of the two
expressions for $L$ into nonzero guarded factors times
\[
 (us-w)G,\qquad G=Kw^2z(u-z)+(u-w)(w-1)(z-1)^2.
\]
On $G=0$, one obtains
\[
 K=\frac{(u-w)(w-1)(z-1)^2}{w^2z(z-u)},\qquad
 L=\frac{us-2wz+w}{w^2z(z-u)}.
\]
The leading coefficient of $P_2$ then equals $L$, an impermissible
extra agreement at infinity.  Thus $u=w/s$, $v=z/s$, where $s\ne0,1$.

For a pair, divide its difference by its two prescribed quadruple-root
factors and take the determinant of the remaining linear forms on each
triple.  Set
\begin{align*}
 A&=Kwz(w+z),& H&=(w-1)(z-1)(s-1),\\
 F_3&=-Kw^2z+(w-1)(z-1)^2,& F_4&=Kwz^2-(w-1)^2(z-1).
\end{align*}
After (G2), the determinants for $245,267,346,357$ are respectively
\[
\begin{array}{cc}
 \displaystyle\frac{2(z-w)(A+H)}{w^2z^2(w-1)(z-1)},&
 \displaystyle\frac{2K(z-w)(A+H)}{wz(w-1)^2(z-1)^2},\\[1ex]
 \displaystyle\frac{2s^2F_3(A-H)}{\Delta},&
 \displaystyle\frac{-2s^2F_4(A-H)}{\Delta},
\end{array}
\qquad \Delta=w^2z^2(w-1)^2(z-1)^2.
\]
All denominators are guarded and $H\ne0$.  Hence $A=-H$ and
$F_3=F_4=0$.  The latter imply $(w-z)(w+z-1)=(w-z)s=0$, again impossible.
The displayed determinant and resultant identities in these three cases
are ordinary polynomial identities, independently checked in the
corresponding \texttt{ORBIT5} and \texttt{ORBIT6} certificate files.

\smallskip\noindent\textbf{Type 7: a complete normalized classification.}
Normalize its quadruple coordinates to $\infty,0,1,u,v,w,z$, subtract
$P_1$, and use amplitudes $j,k,\ell$ for the nonzero values at
$0,1,\infty$.  Put
\[
\begin{array}{lll}
 Q_2=(X-w)(X-z),&Q_3=(X-v)(X-z),&Q_4=(X-u)(X-v),\\
 Q_5=(X-u)(X-z),&Q_6=(X-u)(X-w),&Q_7=(X-v)(X-w).
\end{array}
\]
Its complete quadruple-compatible family is
\begin{align*}
 P_i&=Q_i[(k/Q_i(1)-j/Q_i(0))X+j/Q_i(0)] &&(i=2,4),\\
 P_i&=Q_i[\ell(X-1)+k/Q_i(1)] &&(i=3,6),\\
 P_i&=Q_i[\ell X+j/Q_i(0)] &&(i=5,7).
\end{align*}
Every denominator is a node guard.
Dividing pair differences by their two quadruple-root factors
leaves linear forms.  On each of the seven triples the determinant of
two such forms vanishes.  In the order
$123,145,167,246,257,347,356$, divide these determinants by
\[
 -k,\ -j,\ -\ell,\ -P_2(u),\ [X^3]P_2-\ell,
 \ j-P_3(0),\ -P_3(u),                         \tag{G3}
\]
respectively.  Each divisor is a nonzero outsider--quadruple difference.
The resulting seven expressions are linear in $j,k,\ell$, giving a
$7$-by-$3$ matrix $M$ whose entries are explicit rational functions of
$u,v,w,z$.  Clear only node-guard denominators.

Number rows from zero.  Minor $(0,1,3)$ factors, up to nonzero guards, as
\[
 [(w-1)(z-v)u+vw(w-v)]
 [(v(w+z-1)-wz)u+wz(1-v)].                    \tag{G4}
\]
On the second branch its denominator is nonzero, since otherwise
$wz(1-v)=0$.  The substituted minors $(0,1,2)$ and $(0,3,4)$ are, up to
nonzero guards, the squares of
\[
 v(w+z-1)-2wz+z,\qquad v(w+z-1)-z.
\]
Their simultaneous vanishing contradicts $2z(1-w)\ne0$.  Every
admissible realization therefore satisfies
\[
 u=\frac{vw(v-w)}{(w-1)(z-v)}.                \tag{G5}
\]
This assertion also holds after each relabeling preserving the ordered
pair $(T,C)$ and renormalization of its first three quadruple coordinates.
There are twenty-one such relabelings; only validity of the listed
relabelings, not completeness of that group list, is needed below.

Here is the exact identity certificate for the remaining elimination.
Substitute (G5) in all $35$ minors of $M$ and in the relabeled versions
of (G5).  Cancel only nonzero node guards.  Adjoin a variable $h$ and
the equation $h(v-z)(w-1)-1=0$.  Rational polynomial linear combinations
of these necessary equations give
\[
 (z-1)(w^2+w-z)=0,\qquad
 (z-1)(wz-2w+z-1)=0,\qquad
 (z-1)(v-z+1)=0.                              \tag{G6}
\]
For reproducibility, \texttt{orbit7\_generic.py} constructs $M$ by the
operations just specified, \texttt{orbit7\_branches.py} performs (G4)--(G5),
and \texttt{orbit7\_symmetry.py} checks each relabeling and projective
renormalization.  The identity file
\texttt{orbit7\_identity\_certificate.json.gz} contains $380$ nodes.
An input node is one of these necessary equations.  Every subsequent
node is an explicit sum of polynomial multiples of earlier nodes.
An independent verifier, \texttt{orbit7\_certificate\_independent.py},
checks the inputs against the necessary-equation files and checks all
$359$ subsequent identities by sparse rational arithmetic, then checks
the three output identities (G6).  It uses no symbolic-algebra library.
Thus correctness or completeness of a Groebner-basis implementation is
not a trusted step of this certificate.

Since $z\ne1$, equations (G5)--(G6) imply
\[
 w^3+2w^2-w-1=0,\qquad u=w^2,\quad v=w^2+w-1,\quad z=w^2+w.
\]
On this cubic locus the minor of rows $0,1$ and columns $j,k$ of $M$
is $6w^2-2w-2\ne0$.  The matrix has rank two, and after scaling $j=1$
its unique kernel vector is
\[
 (j,k,\ell)=(1,3w^2+4w-5,3w^2+4w-6).
\]
Nonvanishing follows from irreducibility of the cubic; substitution
checks the kernel vector in all seven rows.  Thus every admissible
type-7 realization has the displayed coefficients, up to the invertible
transformations already made.  Its first five triple coordinates are
\[
 a_1=1-w^2,\quad a_2=w^2+2w-1,\quad a_3=w^2+2w+1,
 \quad a_4=w^2-1,\quad a_5=-w^2-2w+3.
\]
An eighth cubic must match every triple coordinate.  The first three
received values are zero, so it must be
$\lambda\prod_{i=1}^3(X-a_i)$.  The fourth value determines
$\lambda=(1-4w-2w^2)/4$.  At the fifth its difference from the received
value is $486-390w-270w^2$, which is nonzero modulo
$w^3+2w^2-w-1$.  These are direct substitutions in the displayed
polynomials.  This excludes the last possibility in (G1) and proves the
characteristic-zero upper bound.

Finally the argument uses finitely many rational polynomial identities.
Exclude primes dividing their coefficient denominators and the nonzero
integer resultants separating $f(T)=T^3+2T^2-T-1$ from the rank minor and the
interpolation obstruction.  Outside this finite set, and characteristic
two, every step remains valid over an algebraic closure.  Taking $p_0$
larger than those finitely many primes proves the asserted characteristic
range; no optimal value of $p_0$ is claimed.
\end{proof}
